\documentclass[12pt,parskip=half-, twoside=semi]{scrartcl}
\usepackage[T1]{fontenc} % encodage des caractères
\usepackage[french]{babel} % texte en français
\usepackage{microtype} % meilleur justification du texte
\usepackage{cmbright} % police sans serif
\usepackage[frenchmath]{mathastext} % pour le lettres majuscules droites
\usepackage{stmaryrd} % \sslash pour le symbole parallèle
\usepackage{mathtools}
\usepackage{graphicx}
\usepackage{tikz, tkz-euclide}
\usepackage{gensymb}
\usepackage{enumitem}
\setlist[enumerate]{label=\alph*)}
\usepackage{tcolorbox}
\usepackage{gensymb}
\usepackage{multicol}
\NewDocumentCommand{\dotsforanswer}{m o O{1ex}}%
{%
    \IfNoValueTF{#2}%
    {
        \foreach \i in {1,...,#1}{\dotfill\vskip #3}
    }
    {
        #2\dotfill\vskip #3
        \foreach \i in {2,...,#1}{\dotfill\vskip #3}
    }

}
\SetEnumitemKey{cols}
{
    itemsep= 1em,
    before = \raggedcolumns\begin{multicols}{#1},
    after = \end{multicols},
}
\usepackage{tcolorbox}
\newtcbox{\blackbox}{%
    on line,
    left=0pt,
    right=0pt,
    bottom=0pt,
    top=0pt,
    sharp corners,
    colframe=black,
    colback=black,
    coltext=white,
    fontupper=\bfseries
}


%%%%%%%%
% XSIM %
%%%%%%%%
\usepackage[use-files, use-aux]{xsim}
% Pour le template de l'interrogation
\DeclareExerciseEnvironmentTemplate{interrogation_template}%
{%
    \GetExercisePropertyTF{points}%
    % s'il y a des points
    {
        \blackbox%
        {%
            \XSIMmixedcase{\GetExerciseName}~\GetExerciseProperty{counter}%
        }%
        \hfill
        \blackbox{
            \GetExerciseProperty{points}
        }
        \par
    }
    % sinon
    {
        \blackbox%
        {%
            \XSIMmixedcase{\GetExerciseName}~\GetExerciseProperty{counter}%
        }%
        \par
    }%
}%
{%
    \par%
}
\xsimsetup{
    exercise/template = interrogation_template,
    solution/template = interrogation_template,
    exercise/name = question,
}

% Changer le titre ici
\newcommand{\montitre}{Devoir - Mars 2023}
% Changer la date ici
%\newcommand{\madate}{Lundi jj mm aaaa}


% En-tête et pied de page
\usepackage{scrlayer-scrpage}
\setlength{\headheight}{50pt}
\lohead{Nom :\\Prénom :\\Classe :}
\rohead{\no~\fbox{\phantom{XX}}\\~\\ Date : }

\begin{document}
% Titre
\begin{huge}
    \bfseries
    \montitre
    \hfill
    \dots/\TotalExerciseGoal{points}{}{}
    \bigbreak
\end{huge}

\emph{Sur feuille à entête}

\begin{exercise}[points=8]
	Effectue. 
	\begin{enumerate}[cols=2, itemsep=.2cm]
		\item \(4+5\cdot 7-13\)
		\item \(\left(-3+10\right)\cdot 2^3\)
		\item \(\left(8+(-3)\right)^2 - 120\)
		\item \(-4+(-2)-(-7)\)
		\item \(1^9 + 3^3 - 30 + 13\)
		\item \(-10+(-5) + \left(7-13\right)\)
		\item \(45:9\cdot 3 \)
		\item \(-9-13+2\cdot (-5+7)\)
	\end{enumerate}	
\end{exercise}

\begin{exercise}[points=6]
	Complète avec le vocabulaire adéquat : diviseur, multiple.
	\begin{enumerate}[cols=3, itemsep=.5cm]
		\item \(33\) est un  \dotfill \(3\)
		\item \(5\) est un  \dotfill \(15\)
		\item \(0\) est un  \dotfill \(7\)
		\item \(4\) est un  \dotfill \(1\)
		\item \(1\) est un  \dotfill \(9\)
		\item \(8\) est un  \dotfill \(8\)
	\end{enumerate}
\end{exercise}

\begin{exercise}[points=4]
    Trouve le PGCD et le PPCM des nombres suivants. Laisse ta démarche visible.
    \begin{enumerate}[cols=4]
        \item $210$ et $300$
        \item $7 $ et $ 15$
        \item $350$ et $700$
        \item $15$ et $144$
    \end{enumerate}
\end{exercise}

\begin{exercise}[points=8]
    Réduis au maximum. 
    \begin{enumerate}[cols=2]
        \item $3a + 2b - 50a$
        \item $2(10x+3) - (24x-2)$
        \item $3a \cdot (-5a^2b)$
        \item $-2xy^3 \cdot 5x^3 - 2x^2 \cdot 7x^2y^3$
        \item $5x+(2x-3)$
        \item $(5a+2b)-(-3b+8a)$
        \item $5a(a-6) + 2a^2$
        \item $3x+5y-10xy$
    \end{enumerate}
\end{exercise}

\begin{exercise}[points=4]
    Marlène possède une piscine divisée en trois bassins. Le fond du premier est peint en bleu et représente $\dfrac{1}{10}$ de sa piscine. Le fond du deuxième est peint en rose et représente $\dfrac{2}{7}$ de sa piscine. Le troisième bassin a un fond multicolore et représente $\dfrac{2}{6}$ de ce qu'il reste. 
    \begin{enumerate}
        \item Quelle fraction de la piscine de Marlène est peinte en multicolore? 
        \item Quelle fraction de la piscine de Marlène reste-t-il (et pourra être peinte en vert)? 
    \end{enumerate}
\end{exercise}

\begin{exercise}[points=8]
    Effectue. Laisse tes étapes visibles.
    \begin{enumerate}[cols=2]
        \item $4+\dfrac{3}{4}$
        \item $\dfrac{-15}{13} \cdot \dfrac{-39}{45}$
        \item $\dfrac{5}{7} - \dfrac{18}{15}$
        \item $-\dfrac{-27}{45} - 0,5$
        \item $\left(\dfrac{3}{8} + \dfrac{-7}{2}\right)\cdot \dfrac{15}{12}$
        \item $\dfrac{1}{3}-8+\dfrac{-4}{16}$
        \item $\dfrac{-9}{54} + \dfrac{-72}{8} \cdot \dfrac{-32}{27}$
        \item $1,6 + \dfrac{3}{5} - \left(\dfrac{5}{15}\right)^3$
    \end{enumerate}
\end{exercise}

\begin{exercise}[points=3]
	Que doit être la valeur de $\bullet$ pour que 358$\bullet$ soit un nombre divisible par 12? Laisse tes étapes visibles.
\end{exercise}

\begin{exercise}[points=2]
	5 élèves dans la classe de Marion (24 élèves) viennent en marchant à l'école. Quel pourcentage cela représente-t-il? Laisse ta démarche visible et arrondis au dixième près si besoin.
\end{exercise}

\begin{exercise}[points=2]
    Réduis au maximum. Laisse tes étapes visibles.
    \begin{enumerate}[cols=2]
        \item $\dfrac{-2x^3y^5}{8x^2y^9} \cdot \dfrac{-16x^3y}{7}$
        \item $\dfrac{3}{5} - \dfrac{a+2}{4}$
    \end{enumerate}
\end{exercise}

\end{document}